\scanpage{56}
\mspage{27}
Let us return to the case of a general MD-graph, $\underline G$ its modular dimension is
\[
\text{(67)}\qquad\delta(G)=\sum_{\alpha\in S}\underbrace{\bigl(3(g_\alpha-1)+\nu_\alpha+\widetilde\mu_\alpha\bigr)}_{\delta_\alpha}
\]
if we want it to be equal to $1$ or $2$, this means that we are in one of the two cases:
\[
\left\{\begin{array}{l}
1\text{ or }2\text{ of the }\delta_\alpha\text{ are equal to }1,\text{ the others are zero}\\
1\text{ of the }\delta_\alpha\text{ is equal to }2,\text{ the others are zero}
\end{array}\right.
\]
The case $\delta(G)=1$ corresponds to the case where \unsure{exactly} one of the $\delta_\alpha$ is $=1$, the others being zero.

II) $\underline G$ \uline{of modular dimension} $1$

This means that there exists a vertex $\alpha_0$ (and only one, \deleted{in} fact) of $\underline G$ which \deleted{is of genus} \inserted{satisfies the} condition
\[
\begin{aligned}
\text{(68)}\quad &(\alpha_0\text{ of genus }1,\text{ and its total weight }\nu_{\alpha_0}+\widetilde\mu_{\alpha_0}\text{ is equal to }\unsure{1}\\
&\qquad\text{i.e. }\alpha_0\text{ is of order }1,\text{ or isolated \unsure{vertex|point} of marked weight }1)\\
&\uline{\text{or}}\quad(\alpha_0\text{ of genus }0,\text{ and its total weight }\nu_{\alpha_0}+\widetilde\mu_{\alpha_0}=4,\\
&\qquad\text{i.e. }(\nu_{\alpha_0},\widetilde\mu_{\alpha_0})\in\{(4,0),(3,1),(2,2),(1,3),(0,4)\})
\end{aligned}
\]
with in addition
\[
\begin{aligned}
\text{(69)}\quad&\text{the vertices }\alpha\ne\alpha_0\text{ are of genus }0,\\
&\text{of order }\widetilde\mu_\alpha\leq3,\text{ and for such an }\alpha\text{ the marked}\\
&\text{weight }\nu_\alpha\text{ is equal to }3-\widetilde\mu_\alpha.
\end{aligned}
\]
\deleted{We see that the MD-structure of such a graph}
\margin{55}
